\documentclass[11pt]{article}
\usepackage[a4paper,margin=27mm]{geometry}
\usepackage{amsmath,amssymb,amsthm,mathtools,bm}
\usepackage[T1]{fontenc}
\usepackage{lmodern,microtype,xcolor}
\usepackage[colorlinks=true,linkcolor=blue!45!black,citecolor=blue!45!black,urlcolor=blue!45!black]{hyperref}
\numberwithin{equation}{section}
\newtheorem{theorem}{Theorem}[section]
\newtheorem{lemma}[theorem]{Lemma}
\newtheorem{proposition}[theorem]{Proposition}
\newtheorem{corollary}[theorem]{Corollary}
\theoremstyle{remark}\newtheorem{remark}[theorem]{Remark}
\DeclareMathOperator{\Pf}{Pf}
\DeclareMathOperator{\dist}{dist}
\DeclareMathOperator{\supp}{supp}
\newcommand{\ii}{\mathrm i}
\newcommand{\dd}{\mathrm d}
\newcommand{\Rea}{\operatorname{Re}}
\newcommand{\Ima}{\operatorname{Im}}
\newcommand{\T}{\mathbb T}
\newcommand{\X}{\mathcal X_H}
\allowdisplaybreaks[2]
\title{The complete high-temperature Ising susceptibility:\\ a constrained-coordinate proof of the natural boundary}
\author{Dehao Lin\thanks{\raggedright School of Physics, Sun Yat-sen University.\\ORCID: \href{https://orcid.org/0009-0001-4551-8490}{0009-0001-4551-8490}. Contact: \href{mailto:lindh9@mail2.sysu.edu.cn}{\texttt{lindh9@mail2.sysu.edu.cn}}.}}
\date{}
\hypersetup{pdfauthor={Dehao Lin},pdftitle={The complete high-temperature Ising susceptibility: a constrained-coordinate proof of the natural boundary}}
\begin{document}\maketitle
\begin{abstract}
We prove that the complete zero-field isotropic square-lattice Ising
susceptibility, continued from high temperature, has the unit circle
as a natural boundary in $s=\sinh(2\beta_{\rm phys}J)$.
An elementary bordered determinant represents the odd-particle
pair product with a constant final row and column. On the exact
full-lattice angular constraint a balanced deformation preserves the
total angle. A regular coordinate and one collective tangential field
then give a logarithmic absolute bound for the entire differentiated
higher odd-particle sum. The same estimate controls finite-order
complements. At a dense family of odd prime roots, two reflected
collision charts give a nonzero inverse-distance derivative, evaluated
by a rational radial primitive. Signed constraints treat the two
boundary halves, with the physical normalization and the origin
included throughout.
\end{abstract}

\section{The physical function and the proof}
The high-temperature susceptibility of the square-lattice Ising model
is represented by an infinite series of odd-particle integrals. The
singularities of individual terms accumulate densely on the unit circle,
and their survival in the complete sum is the central analytic issue.
Orrick, Nickel, Guttmann, and Perk \cite[Section~3.2]{Orrick} and
Boukraa et al.\ \cite{Boukraa} studied the singularities and the series
evidence for a natural boundary. We use a balanced deformation of the constrained angular integral,
a regular inverse coordinate, and a single tangential field to
estimate the full differentiated tail.

Let $J>0$ and consider the zero-field isotropic Ising model on the
infinite square lattice. For $0<s=\sinh(2\beta_{\rm phys}J)<1$,
let
\[
 \X(s)=\beta_{\rm phys}^{-1}\chi_H(s)
       =\sum_{(m,n)\in\mathbb Z^2}\langle\sigma_{00}\sigma_{mn}\rangle_s.
\]
The high-temperature infinite-volume Gibbs state has zero
magnetization. This sum includes every lattice site and its origin.
We use $t=s^{-1}$ only as an exterior coordinate and put
\[
 S(t)=t+t^{-1},\qquad Q_H(t)=t(1-t^{-4})^{1/4},
\]
where the fourth root tends to one at infinity.

\begin{theorem}\label{thm:main}
The physical high-temperature germ of $\chi_H$ admits no local
holomorphic continuation across any point of $|s|=1$.
The same is true of the normalized complete susceptibility $\X$.
\end{theorem}

The real-temperature physical inputs are the odd form-factor
expansion of the full correlation and its full-lattice angular
sum \cite[eqs.~(1.9), (1.11), (3.13), (3.17)--(3.20),
(4.1)--(4.2)]{Orrick}, also recorded in
\cite[eqs.~(3), (5)--(9)]{Boukraa}. Their use is restricted to the
physical real interval. We establish complex continuation and the
boundary estimates below.
For the complete angular coefficients $C_N$ we prove
\[
 \X(1/t)=Q_H(t)\sum_{\substack{N\ge1\\N\ {\rm odd}}}C_N(t).
\]
At $t_*=e^{\ii\alpha}$, $\alpha=2\pi a/p\in(0,\pi)$, for each
odd prime $p\ge5$, set $k=(p^2-1)/2$. The decisive estimates are
\[
 C_p^{(k)}((1+\varepsilon)t_*)
       =\Lambda_\alpha\varepsilon^{-1}+o(\varepsilon^{-1}),
 \qquad \Lambda_\alpha\ne0,
 \qquad
 \sum_{\substack{N>p\\N\ {\rm odd}}}|C_N^{(k)}((1+\varepsilon)t_*)|
       =O(1+|\log\varepsilon|).
\]
The needed fixed derivatives of lower particle orders are logarithmically bounded.
Thus the complete infinite sum cannot be holomorphic at these dense
points. Dividing by the locally nonzero physical factors transfers
the obstruction to $\chi_H$.

\subsection*{AI Declaration}
ChatGPT was used extensively in the mathematical research, proof development,
and preparation of this manuscript.

\section{The exact full-lattice angular representation}
\label{sec:representation}\label{sec:angular-bulk}
Write
\[
 W(y;t)=S(t)-\tfrac12(y+y^{-1}),\qquad
 z=W-\sqrt{W^2-1},\qquad
 \sinh\gamma=\tfrac12(z^{-1}-z),\quad z=e^{-\gamma}.
\]
For $|t|>1$ and $|y|=1$ choose the unique root with $|z|<1$,
continued from the positive real interval. Indeed $S(t)\notin[-2,2]$,
so $W$ cannot lie in $[-1,1]$. The root and $\sinh\gamma$ are
single-valued there.
For $y_i=e^{\ii\theta_i}$ define the complete pair
\begin{equation}\label{eq:kernel-square}
 P_{ij}=
 \left[\frac{\sin((\theta_i-\theta_j)/2)}
 {\sinh((\gamma_i+\gamma_j)/2)}\right]^2
 =-\frac{(y_i-y_j)^2z_i z_j}{y_i y_j(1-z_i z_j)^2}.
\end{equation}
Only the rational square is used globally, so logarithmic and
half-angle sign choices do not affect the integral.
Let $Z=\prod_i z_i$, and for $\sigma\in\{1,-1\}$ set
$a_+=0$, $a_-=\pi$, and
\[
 \mathbb H_N^\sigma=\{\theta\in\T^N:\textstyle\sum_i\theta_i
                       =a_\sigma\pmod {2\pi}\},\qquad
 d\mu_N=\prod_{i<N}\frac{\dd\theta_i}{2\pi}.
\]
\[
 \theta_N=a_\sigma-\sum_{i<N}\theta_i\pmod {2\pi}.
\]
This is normalized Haar measure; its Euclidean hypersurface
normalization is $[\sqrt N(2\pi)^{N-1}]^{-1}$.
For odd $N$ define
\begin{equation}\label{eq:angular-bulk}
 C_N^\sigma(t)=\frac1{N!}\int_{\mathbb H_N^\sigma}
    \frac{1+\sigma Z}{1-\sigma Z}
    \prod_i\frac1{\sinh\gamma_i}\prod_{i<j}P_{ij}\,d\mu_N,
 \qquad C_N=C_N^+.
\end{equation}
The minus sign is an auxiliary angular constraint, not a different
physical observable in the theorem.

\begin{proposition}[Physical identification and normal convergence]
\label{prop:angular-bulk}
The complete normalized physical susceptibility has the exterior
continuation
\begin{equation}\label{eq:bulk}
 \X(1/t)=Q_H(t)\sum_{\substack{N\ge1\\N\ {\rm odd}}} C_N(t).
\end{equation}
Both signed series, with each fixed number of derivatives, converge
normally on compact subsets of $|t|>1$.
The origin contributes exactly one to $\X$.
\end{proposition}
\begin{proof}
For real $t>1$, the per-site physical input is
\[
 \langle\sigma_{00}\sigma_{mn}\rangle_{1/t}
  =Q_H(t)\sum_{\substack{N\ge1\\N\ {\rm odd}}}c_N(m,n;t),
\]
\[
 c_N(m,n;t)=\frac1{N!}\int_{\T^N}
 Z^{|m|}e^{\ii n\sum\theta_i}\prod_{i<j}P_{ij}
                   \prod_i\frac{\dd\theta_i}{2\pi\sinh\gamma_i}.
\]
The exponential is equivalent under angle reversal to the cosine
in \cite[eq.~(4.2)]{Orrick}. In the isotropic specialization
$\cot\alpha_{\rm anis}=1$, its $k=s^2=t^{-2}$ and
$k^{-1/2}(1-k^2)^{1/4}=Q_H(t)$; the symbols $\alpha_{\rm anis}$
and $k$ here belong only to that source normalization.
The dispersion and pair factors agree literally with
\eqref{eq:kernel-square}.

The bordered determinant identity proved in
Section~\ref{sec:elementary-pfaffian} has size $N+1$ with bounded
entries on a common small complex angular strip. There the
one-body factors are bounded and $|z_i|\le e^{-c}<1$.
Shifting every angle by $\ii\eta\operatorname{sgn}(n)$, for
fixed sufficiently small $\eta>0$, gives
\[
 |c_N(m,n;t)|\le \frac{C^N N^{(N+1)/2}}{N!}
                   e^{-N(c|m|+\eta|n|)}.
\]
Periodicity justifies the shift; Hadamard supplies the determinant
bound. This is summable in $m,n$ and all odd $N$.
Thus particle and lattice sums may be interchanged.
The $m$-sum gives $(1+Z)/(1-Z)$.
For a periodic analytic $F$ on $\T^N$, define
\[
 A(v)=\int_{\T^{N-1}}
 F(\theta_1,\ldots,\theta_{N-1},v-\textstyle\sum_{i<N}\theta_i)
                   \prod_{i<N}\frac{\dd\theta_i}{2\pi}.
\]
Its Fourier coefficient of index $-n$ is the $N$-angle integral
with multiplier $e^{\ii n\sum\theta_i}$.
The absolutely convergent Fourier sum is $A(0)$, giving precisely
\eqref{eq:angular-bulk} with the normalized Haar measure $d\mu_N$.
At the origin the per-site sum is $\langle\sigma_{00}^2\rangle=1$,
so the origin is already contained in this full-lattice sum.
For a further normalization check,
\[
 C_1(t)=\frac1{S(t)-2}=\frac{s}{(1-s)^2}.
\]

On a compact exterior neighborhood and a slightly larger one,
$m\le|z_i|\le q<1$ uniformly on the angular torus.
All one-body and bordered entries are bounded, and the signed
product denominator has a fixed gap. Since the Haar measure is one,
\[
 |C_N^\sigma(t)|\le C_0\sqrt N(C_1/\sqrt N)^N.
\]
The bound is summable; Cauchy's estimate on the larger compact
neighborhood proves normal convergence of every fixed derivative.
The identity theorem gives \eqref{eq:bulk} on the exterior.
Near $s=0$, the interior root has the form $z=s\,a(s,y)$ with
analytic bounded $a$, and $1/\sinh\gamma=s\,b(s,y)$ with
analytic bounded $b$. The same determinant estimate gives normal
convergence of $\sum C_N(1/s)/s$. Multiplication by
$(1-s^4)^{1/4}$ therefore fills the removable point $s=0$ and
recovers the full high-temperature analytic germ.
\end{proof}

The signed coefficients commute with complex conjugation and obey,
for odd $N$,
\begin{equation}\label{eq:negation}
 C_N^+(-t)=-C_N^-(t).
\end{equation}
Indeed replace $y_i,z_i,t$ by $-y_i,-z_i,-t$.
The angular constraint becomes $\prod y_i=-1$, every pair is
unchanged, all $N$ one-body factors change sign, and $Z$ changes
sign. Angular translation preserves Haar measure.
In particular no evenness of the physical high-temperature
susceptibility is asserted.

\section{An equal-angle odd-prime family}
\label{sec:prime-roots}
\begin{lemma}\label{lem:prime-root-family}
Let $p\ge5$ be prime, $1\le a\le(p-1)/2$, and put
\[
 \alpha=2\pi a/p,\quad t_*=e^{\ii\alpha},\quad
 S_*=2\cos\alpha,\quad N_*=p,\quad k=(p^2-1)/2.
\]
These points are dense in both open upper quadrants, $S_*\ne0$,
and $k\equiv0\pmod4$.
The first odd order admitting
$S_*=\cos(2\pi l/n)+\cos(2\pi m/n)$ is $p$.
At that order the unordered cosine pair is uniquely
$\{\cos\alpha,\cos\alpha\}$.
For $b=\arccos(|S_*|-1)$, the algebraic number $e^{-\ii b}$
is not a root of unity.
\end{lemma}
\begin{proof}
Write $d=(p-1)/2$, $\zeta_p=e^{2\pi\ii/p}$,
$u_j=\zeta_p^j+\zeta_p^{-j}$, and
$E=\mathbb Q(\zeta_p+\zeta_p^{-1})$.
The $u_1,\ldots,u_d$ are linearly independent over $\mathbb Q$:
a relation gives a polynomial of degree at most $p-1$, with zero
constant term, vanishing at $\zeta_p$; it is a multiple of
$1+X+\cdots+X^{p-1}$, and that multiple must be zero.
The conjugates of each $u_a$ are exactly $u_1,\ldots,u_d$,
with $u_1>0$ and $u_d<0$.
Also $u_a$ is irrational: otherwise $\zeta_p^a$ would satisfy
$X^2-u_aX+1$ over $\mathbb Q$, contrary to its degree $p-1>2$.
If $p\nmid n$, the cyclotomic fields
$\mathbb Q(\zeta_p)$ and $\mathbb Q(\zeta_n)$ have intersection
$\mathbb Q$: their compositum has degree
$\varphi(pn)=\varphi(p)\varphi(n)$.
A cosine representation at an odd $n<p$ would contradict
irrationality. Order $p$ realizes the equal pair.

Put $\delta=\operatorname{sgn}(u_a)$ and $S_+=\delta u_a=|S_*|$.
The number $S_+$ has a negative conjugate, using $u_d$ when
$\delta=1$ and $u_1$ when $\delta=-1$.
If $\xi=e^{-\ii\arccos(S_+-1)}$ were a root of unity,
every conjugate of $\xi+\xi^{-1}=2S_+-2$ would lie in $[-2,2]$.
The indicated conjugate is strictly less than $-2$, a contradiction.
Here each embedding of $\mathbb Q(S_+)$ extends to the field
containing $\xi$. This proves the asserted endpoint nonresonance
for both signs at once.

At odd order $p$ the endpoint cosine $-1$ is absent.
A cosine $1$ has an impossible partner $S_*-1<-1$ when
$S_*<0$; for $S_*>0$ it is excluded by the preceding endpoint
argument. All remaining twice-cosines are $u_j$.
The equation $2u_a=u_l+u_m$ and independence force $l=m=a$
modulo reflection. The only angular equality centers are therefore
$\theta_i=-\alpha$ and $\theta_i=+\alpha$ for every $i$.
Approximation of any $\alpha_0\in(0,\pi)$ by nearest prime-grid
angles proves density. Finally every odd square is $1$ modulo $8$.
\end{proof}

For a point with negative $S_*$, let
$\widehat t_*=-\overline{t_*}=e^{\ii(\pi-\alpha)}$ and use
$C_N^-$ there. Equation~\eqref{eq:negation} gives
\begin{equation}\label{eq:negative-derivatives}
 C_N^{(j)}((1+\varepsilon)t_*)
   =(-1)^{j+1}\overline{(C_N^-)^{(j)}((1+\varepsilon)\widehat t_*)}.
\end{equation}
A signed equality at odd order $n$ has both angles in
$(2\pi\mathbb Z+\pi)/n$. Adding $\pi$ to both angles
converts it to an ordinary odd-order equality at $-S_*$.
Thus the same arithmetic proves uniqueness and absence of lower
orders on the signed half. Nonresonance gives
$e^{-\ii nb}\ne\pm1$ for every integer $n\ge1$.

\section{A bordered determinant for odd particles}
\label{sec:elementary-pfaffian}\label{sec:determinant-local}
The physical pair product has a determinant with a constant border.
This keeps both its collision zeros and its factorial estimates visible.
\begin{lemma}[Four points on the dispersion curve]
Let nonzero $y_i,z_i$ satisfy
$y_i+y_i^{-1}+z_i+z_i^{-1}=2S$ for $1\le i\le4$, and set
$B_{ij}=(y_i-y_j)/(1-z_i z_j)$.  Then
\[
 B_{12}B_{34}-B_{13}B_{24}+B_{14}B_{23}
 =\left(\prod_{i=1}^4\frac{z_i}{y_i}\right)
   \prod_{i<j}B_{ij}.
\]
The equality holds wherever the expressions are defined and,
after clearing denominators, at every point of the curve.
\end{lemma}
\begin{proof}
Write $a_{ij}=y_i-y_j$, $t_i=z_i+z_i^{-1}$, and
$Z_4=z_1z_2z_3z_4$.  For a partition $ij\mid kl$ of the four
indices, direct multiplication gives
\[
 \frac{\prod_{v\in\{i,j\},\,w\in\{k,l\}}(1-z_vz_w)}{Z_4}
 =R+t_it_j+t_kt_l,
 \qquad
 R=Z_4+Z_4^{-1}+2-\sum_{v<w}
          \bigl(z_vz_w+(z_vz_w)^{-1}\bigr).
\]
Multiplying the proposed identity by
$\prod_{i<j}(1-z_i z_j)$, the common $R$ term cancels because
$a_{12}a_{34}-a_{13}a_{24}+a_{14}a_{23}=0$.
The remaining left side, divided by $Z_4$, is
\[
 \begin{split}
 &a_{12}a_{34}(t_1t_2+t_3t_4)
 -a_{13}a_{24}(t_1t_3+t_2t_4)\\
 &\hspace{25mm}+a_{14}a_{23}(t_1t_4+t_2t_3)
 =\det[\,1,\ y_i,\ t_i,\ y_it_i\,]_{i=1}^4.
 \end{split}
\]
Here the last equality is a two-column Laplace expansion.
Since $t_i=2S-y_i-y_i^{-1}$, column operations give
\[
 \det[1,y_i,t_i,y_it_i]
 =\det[1,y_i,y_i^{-1},y_i^2]
 =\frac{\det[y_i,y_i^2,1,y_i^3]}{\prod_i y_i}
 =\frac{\prod_{i<j}(y_i-y_j)}{\prod_i y_i}.
\]
This is the required cleared numerator.
\end{proof}

\begin{lemma}[Pfaffian propagation]\label{lem:pf-propagation}
Suppose a skew matrix $K$ has the property that every four-index
principal submatrix satisfies
\[
 K_{ab}K_{cd}-K_{ac}K_{bd}+K_{ad}K_{bc}
   =\prod_{v<w\in\{a,b,c,d\}}K_{vw}
 \quad(a<b<c<d).
\]
Then, at every even order for which the intervening products
can be canceled, $\operatorname{Pf}K=\prod_{i<j}K_{ij}$.
For a rational or analytic family with a dense set of such
points, the identity extends to the whole family.
\end{lemma}
\begin{proof}
Let $I$ be the first $N-4$ indices and let $a,b,c,d$ be the last
four.  The elementary Pfaffian condensation identity is
\[
 \operatorname{Pf}K_{Iabcd}\operatorname{Pf}K_I
 =\operatorname{Pf}K_{Iab}\operatorname{Pf}K_{Icd}
  -\operatorname{Pf}K_{Iac}\operatorname{Pf}K_{Ibd}
  +\operatorname{Pf}K_{Iad}\operatorname{Pf}K_{Ibc}.
\]
For completeness, when $K_I=A$ is invertible, the simultaneous
row and column elimination in
$\left[\begin{smallmatrix}A&B\\-B^T&C\end{smallmatrix}\right]$
has determinant one and gives the diagonal blocks
$A$ and $C+B^TA^{-1}B$.  Pfaffian multilinearity therefore
reduces the displayed formula to the defining four-by-four
Pfaffian formula.  The identity is polynomial, so it holds
also when $A$ is singular.

Assume the product formula at all smaller even orders.
Writing $P_I=\prod_{i<j\in I}K_{ij}$, the right side factors as
\[
 P_I^2\prod_{i\in I,\,v\in\{a,b,c,d\}}K_{iv}
 \bigl(K_{ab}K_{cd}-K_{ac}K_{bd}+K_{ad}K_{bc}\bigr)
 =P_I\prod_{i<j}K_{ij}.
\]
The left side is $P_I\operatorname{Pf}K$.  Cancel $P_I$.
The orders two and four start the induction.
\end{proof}


\begin{proposition}[Odd bordered identity]\label{prop:complete-product}
\label{prop:local-determinant-coefficient}
For odd $N$, let $\mathcal A$ be the skew $(N+1)$-by-$(N+1)$
matrix
\[
 \mathcal A_{ij}=L_{ij}:=-\ii\frac{y_i-y_j}{1-z_i z_j}
 \quad(i,j\le N),\qquad
 \mathcal A_{i,N+1}=1,\quad \mathcal A_{N+1,i}=-1.
\]
On the dispersion curve $y_i+y_i^{-1}+z_i+z_i^{-1}=2S$,
\begin{equation}\label{eq:odd-border}
 \det\mathcal A=\prod_{i<j}P_{ij}.
\end{equation}
The identity is rational and has no square-root choices.
\end{proposition}
\begin{proof}
Choose local $r_i^2=z_i/y_i$ and put
$K_{ij}=-\ii r_i r_j(y_i-y_j)/(1-z_i z_j)$.
The four-point lemma gives
$\Pf K_{\{1,2,3,4\}}=\prod_{i<j\le4}K_{ij}$:
both sides equal $-(\prod_i r_i^3)\prod_{i<j}B_{ij}$.
Pfaffian propagation proves this product identity at every even
order, first at generic points and then by continuation.

Append one point $y_0=R$ on the curve, using the branch
$z_0=-R^{-1}+O(R^{-2})$ and
$r_0=-\ii R^{-1}+O(R^{-2})$ as $R\to\infty$.
Then $K_{i0}\to r_i$.
Squaring the even identity at order $N+1$ gives
\[
 \det\begin{pmatrix}K&r\\-r^T&0\end{pmatrix}
       =\left(\prod_i r_i^2\right)\prod_{i<j}K_{ij}^2.
\]
Apply the congruence $A\mapsto DAD$ with
$\operatorname{diag}(r_1^{-1},\ldots,r_N^{-1},1)$.
Its determinant is multiplied by $\prod_i r_i^{-2}$,
and the resulting matrix is exactly $\mathcal A$.
Since $K_{ij}^2=P_{ij}$, this proves \eqref{eq:odd-border}.
The denominators may be cleared to extend across removable cases.
For $N=1$, the matrix is $\left[\begin{smallmatrix}0&1\\-1&0\end{smallmatrix}\right]$
and the identity reads $1=1$.
\end{proof}

\begin{lemma}[Minors and collisions]\label{lem:border-minors}
If the ordinary entries $L_{ij}$ are bounded by $a$, a determinant
with $\ell$ marked ordinary entries, including any fixed derivatives
of them, has its unmarked factor bounded by
\[
 C_j N^{C_j}(C\sqrt N\,a)^{N-1-\ell}
\]
when the marked entries and their fixed jets are bounded.
Here $\ell\le j$ and $j$ is fixed; terms with more marked entries
than the polynomial permits vanish.
On a regular fixed-order collision chart, $\det\mathcal A$
has the complete squared Vandermonde factor of degree $N(N-1)$.
\end{lemma}
\begin{proof}
A differentiated determinant expands into polynomially many
choices of marked rows and columns, leaving complementary minors.
The ghost row and column are never differentiated because their
entries are constant. Expand both in each remaining minor.
Each nonzero term contains exactly two border entries and a minor
of the ordinary block of size $N-1-\ell$.
There are at most $N^2$ choices; Hadamard bounds that minor.
Fixed repeated derivatives of a marked entry do not remove
additional rows or columns, and their jets have the stated bound.
For the collision assertion use \eqref{eq:odd-border}:
each rational pair has its double zero at $y_i=y_j$.
Permutation covariance of the coupled contour gives the same
zero at $\theta_i=\theta_j$. On a compact regular collision
chart its analytic or smooth quotient has bounded fixed jets.
\end{proof}
\section{A product-preserving deformation of the whole-bulk integral}
\label{sec:constrained-deformation}

Fix either signed family from Section~\ref{sec:prime-roots},
a positive-geometry point $t_*$ with $S_*=t_*+t_*^{-1}\in(0,2)$,
and its sign $\sigma$. Write $\mathcal C_N=C_N^\sigma$,
$N_*=p$, $\widehat\alpha=\arccos(S_*/2)\in(0,\pi/2)$, and
\[
 t_\varepsilon=(1+\varepsilon)t_*,\qquad
 b=\arccos(S_*-1),\qquad \xi=e^{-ib},\qquad
 H=\log(1/\varepsilon).
\]
The constants in this section may depend on the selected point and a
fixed derivative order $k$, but never on the odd particle number $N$
or on $\varepsilon$.  The infinite-tail estimates concern $N\ge N_*+2$, where
$k=(N_*^2-1)/2$.
We prove logarithmic absolute bounds through $k$, which are
smaller than the first singular term. The deformation and the regular transport formulas
are also used at fixed odd orders $3\le N\le N_*$ in
Section~\ref{sec:first-single}.

For the regular dispersion coordinate we write
\[
 \varphi=\arccos W,\qquad z=e^{-i\varphi},\qquad
 \gamma=i\varphi.
\]
Its physical continuation lies in the fourth quadrant near either
inverse-dispersion branch.  The angular integration coordinates remain
$\theta_i$, with $\sum_i\theta_i=a_\sigma\pmod {2\pi}$.

Choose a smooth periodic function $p$ compactly supported in $(0,\pi)$,
with $0\le p\le1$ and $p=1$ on $[\delta,\pi-\delta]$.
Here $\delta>0$ is smaller than $\widehat\alpha/2$, $b/2$, and
$(2\widehat\alpha-b)/2$, and is small enough that $b$ is inside this
plateau and $1+g'(\theta)>0$ on the upper endpoint sliver.
Here $g(\theta)=\arccos(S_*-\cos\theta)$; the choices are
possible because $b<2\widehat\alpha$ and $g'(0)=0$.  Put $m(\theta)=p(-\theta)$.  Thus $p,m$ have
disjoint supports, and both inverse branches have fixed plateau
neighborhoods.  Define
\begin{equation}\label{eq:constrained-homotopy}
 \begin{aligned}
 P&=\sum_i p(\theta_i),\qquad Q=\sum_i m(\theta_i),\\
 h_i&=\frac\tau N\{P m(\theta_i)-Q p(\theta_i)\},\qquad
 y_i(\lambda)=e^{i\theta_i+\lambda h_i},\quad 0\le\lambda\le1.
 \end{aligned}
\end{equation}
The positive number $\tau$ is fixed and sufficiently small.  In
particular $\sum_i h_i=0$, so $\prod_i y_i=\sigma$ throughout the
deformation.

\begin{lemma}[Exact constrained deformation]
\label{lem:constrained-homotopy}
The full integral defining $\mathcal C_N$ is unchanged by
\eqref{eq:constrained-homotopy}.  Its angular Jacobian and its fixed
jets have bounds $C_j^N N^{C_j}$, uniformly in $N$.
\end{lemma}
\begin{proof}
For $y=e^{v+i\theta}$,
\[
 \Im W=\Im S-\sinh v\sin\theta.
\]
On the upper semicircle the motion in
\eqref{eq:constrained-homotopy} is inward, and on the lower
semicircle it is outward.  Both motions increase $\Im W$.
For the admissible exterior parameters it remains positive, so every
$z_i$ stays on the interior sheet.  Thus $1-\sigma Z$ and all the pair
denominators $1-z_i z_l$ stay nonzero.  The kernel square \eqref{eq:kernel-square} has no
$1-y_i y_l$ denominator.  The closed constrained torus therefore
undergoes a legal holomorphic homotopy.

Write $w_i=\theta_i-i h_i$, so $\sum_iw_i=\sum_i\theta_i$.
The full matrix $I-iDh$ preserves the tangent hyperplane and induces
the identity on its one-dimensional quotient.  Its full determinant
is consequently its constrained tangent determinant.  Directly,
\[
 Dh=\frac\tau N\operatorname{diag}(P m'_i-Qp'_i)
       +\frac\tau N(m\otimes p'-p\otimes m').
\]
The diagonal and the two rank-one terms have operator norms at most
$C\tau$, independently of $N$.  Hence the Jacobian has modulus at
most $C^N$, and its inverse matrix is uniformly bounded after reducing
$\tau$.  Differentiating the determinant, or its logarithmic derivative
formula, shows that fixed jets add only a polynomial in $N$.
\end{proof}

Let $\eta$ be smooth, $0\le\eta\le1$, with $\eta(\mu)=0$ for
$\mu\le1/4$ and $\eta(\mu)=1$ for $\mu\ge1/2$.  Partition the final
deformed integral with the three weights
\[
 w_F=\eta(P)\eta(Q),\qquad
 w_L=\tfrac12(1-\eta(P))(1+\eta(Q)),\qquad
 w_U=\tfrac12(1-\eta(Q))(1+\eta(P)).
\]
Their sum is one.  Reflection $\theta\mapsto-\theta$ exchanges $P,Q$,
takes $y_i$ to $y_i^{-1}$, and preserves $z_i$.
Every rational pair in \eqref{eq:kernel-square} is unchanged;
so is the bordered determinant by \eqref{eq:odd-border}.
The constraint $\prod y_i=\sigma$ and Haar measure are preserved
for either sign. Therefore
\begin{equation}\label{eq:constrained-split}
 \mathcal C_N=F_N+2L_N,
\end{equation}
where $F_N$ and $L_N$ are the final integrals with weights $w_F$
and $w_L$.  All weights are imposed after the full closed homotopy.
The support of $L_N$ has $P\le1/2$; it excludes the upper inverse
branch and consists of the lower semicircle together with two small
upper endpoint slivers.

\subsection{The sector containing both branch populations}

\begin{lemma}\label{lem:constrained-selected}
For every fixed $j\ge0$ and every odd $N\ge5$,
\[
 |F_N^{(j)}(t_\varepsilon)|
    \le C_j N^{3j+2}\left(\frac{C_j}{\sqrt N}\right)^N.
\]
In particular these derivatives have a uniformly bounded absolute sum.
\end{lemma}
\begin{proof}
On this support $P,Q\ge1/4$.  Both bump supports stay a fixed
distance from the endpoints $0,\pi$.  The imaginary-part formula,
and the bounded range of $W$, imply
\[
 -\log|Z|\ge c\left(N\varepsilon+\frac{\tau PQ}{N}\right)
             \ge \frac cN.
\]
Every true lower branch has outward displacement $\tau P/N$,
and every true upper branch has inward displacement $\tau Q/N$.
Their normal forms give a branch distance at least $c/N$ in
the $W$ plane.  All other coordinates stay in fixed regular
neighborhoods.  Consequently $|\partial_t\log z_i|\le C\sqrt N$
on a disk $|t-t_\varepsilon|\le c/N^3$.  Reduce $c$ once: the
variation of $\log Z$ is at most $CN^{3/2}c/N^3$, smaller than
half its preceding gap.  The true branches remain on their interior
sheets throughout this disk.

The Cauchy disk may cross $\Im S(t)=0$. All factors on that disk
are understood as analytic continuations of the branches fixed at
$t_\varepsilon$, rather than as newly selected physical roots at each
parameter value. On the true-branch neighborhoods the deformation
margin is of order $1/N$ in the $W$ plane; on the complementary
compact angular sets the roots remain separated from $W=\pm1$.
After reducing the disk constant, these facts give holomorphic
branches for each $z_i$ and the continued bordered kernel.
Moreover,
$|\log Z(t)-\log Z(t_\varepsilon)|
 \le CN^{3/2}|t-t_\varepsilon|$,
so the existing $c/N$ damping gap yields
$|Z(t)|\le\exp(-c'/N)$ and
$|1-\sigma Z(t)|\ge c'/N$ on the entire disk.
The pair-denominator separations below persist there by the same
branch-margin estimates. No sign condition on $\Im S(t)$ is imposed
on the continued integrand.

We estimate the bordered determinant directly.
Entries with both points near the same inverse branch are bounded:
the common factor $\varphi_i+\varphi_l$ cancels, leaving
$(\varphi_i-\varphi_l)$ times a bounded analytic coefficient.
Away from an opposite-branch pair, compactness gives the same bound.
Let $B$ be the number of coordinates in the two fixed true-branch
neighborhoods.  If a pair is near opposite branches, its numerator
is bounded and
\[
 |1-z_i z_l|\ge c\sqrt{(P+Q)/N},\qquad
 |L_{il}|\le C\sqrt{N/(P+Q)}\le C\sqrt{N/B}.
\]
Indeed both regular branch coordinates remain in a common sector
of width strictly less than $\pi$ (a small fixed enlargement of the
fourth quadrant suffices), while their imaginary dispersion margins are
at least $c\tau P/N$ and $c\tau Q/N$.  Each such coordinate
contributes one to $P+Q$.  There are at most $B$ exceptional
entries in any row; all remaining entries are bounded.  Every ordinary row therefore has Euclidean norm at most
$C\sqrt N$; the constant border row has norm $\sqrt N$ as well.  Hadamard gives
\[
 |\det\mathcal A|\le C^N N^{(N+1)/2}
\]
uniformly on the disk.

The one-body factors have a common independent angular majorant
\[
 f(\theta)=C\left(1+
  \operatorname{dist}(\theta,\{b,-b\})^{-1/2}\right).
\]
To verify this, use the imaginary part of the branch normal form
when $|u|\le C(\varepsilon+|h|)$ and its linear real part outside
that interval.  The small parameter disk is absorbed in the positive
branch displacement.  Thus $f\in L^{4/3}(\mathbb T)\cap L^1(\mathbb T)$.
With normalized circle measure, the constrained product integral is
$f^{*N}(a_\sigma)$.  Indeed, $g=f*f$ belongs to $L^2$ by Young's
inequality, and $g*g$ is bounded and continuous by
Cauchy--Schwarz and continuity of translations in $L^2$.
Further convolution with $f\in L^1$ preserves these properties.
Thus point evaluation is legitimate, and Young's inequalities give
\[
 f^{*N}(a_\sigma)\le \|f\|_{4/3}^{4}\|f\|_1^{N-4}\le C^N,
 \qquad N\ge4.
\]
The Jacobian costs $C^N$ by Lemma~\ref{lem:constrained-homotopy}.
The remaining product kernel costs $CN$.  Division by
$N!\ge(N/e)^N$ bounds the integral on the disk by
$N^2(C/\sqrt N)^N$.  Cauchy's inequality proves the assertion.
\end{proof}

\section{A regular coordinate for the one-sided contour}\label{sec:global-coordinate}


The support of $L_N$ can be lifted to one angular interval
$(-\pi-\delta,\delta)$: the remaining upper plateau has zero
weight because $p(\theta_i)=1$ would imply $P>1/2$.
The bordered entries are rational and periodic in each angle,
so this lift preserves the determinant and its integral.
We use $w_i=\theta_i-ih_i$ and the physical phase
$\varphi_i=\arccos(S-\cos w_i)$ throughout.

Choose a fixed branch width, decreased further below if needed, so that
\begin{equation}\label{eq:finite-mean-exclusion}
 \delta_B<\frac14\min_{\substack{3\le n\le p\\n\ \mathrm{odd}}}
       \operatorname{dist}(-b,(2\pi\mathbb Z+a_\sigma)/n).
\end{equation}
The minimum is positive because $e^{-\ii nb}\ne\sigma$ for every $n$.
This finite gap is used only for the lower-order complement in
Section~\ref{sec:first-single}; the infinite-tail estimate below
does not require any mean to avoid the branch.
The right interval will be $[-b-\delta_B,\delta]$; its
complement on the one-sided support has a compact left root.
The displacement strength is chosen after these fixed margins and
cutoffs.

Use the sum coordinate
\[
 v=\varphi+w.
\]
It is regular even at the lower inverse-dispersion branch.

\begin{lemma}[Regular inverse and right-side convexity]\label{lem:inverse}
The one-sided limiting dispersion curve has a fixed complex
neighborhood on which $v$ is a regular coordinate and $w=U(t,v)$
is analytic, with bounded fixed joint parameter and coordinate jets.
On a fixed convex neighborhood $\mathcal D_R$ of the right
$v$-interval, enlarged to contain a short negative branch segment,
\[
 \Rea U_{vv}(t,v)\ge c>0.
\]
The actual right-contour images lie in $\mathcal D_R$ after the
branch width, deformation and parameter neighborhood are fixed small.
\end{lemma}
\begin{proof}
On the limiting compact right arc put
$g(\theta)=\arccos(S_*-\cos\theta)$.
On $[-b,0]$,
$g'(\theta)=-\sin\theta/\sin g(\theta)\ge0$.
Since $g'(0)=0$, a short upper sliver retains
$1+g'(\theta)>c>0$. Thus $v=g(\theta)+\theta$ is strictly increasing.
At the lower branch use $\varphi$ itself:
$w=-\arccos(S(t)-\cos\varphi)$ is even analytic in $\varphi$,
so $\dd(w+\varphi)/\dd\varphi=1$ there. Thus $v$ is regular
through the branch.

There is an explicit right inverse:
\begin{equation}\label{eq:v-1}
 U(t,v)=\frac v2-
       \arccos\left(\frac{S(t)}{2\cos(v/2)}\right).
\end{equation}
Its branch is fixed by $w<\varphi$ on the limiting right arc.
The real $v$-interval is contained strictly in
$(-2\widehat\alpha,2\widehat\alpha)$, since $b<2\widehat\alpha$
and the upper sliver is short. Put $x=v/2$ and
$D=4\cos^2x-S_*^2>0$. Direct differentiation gives
\begin{equation}\label{eq:v-2}
 U_{vv}(t_*,v)=\frac{S_*}{4\cos^2x\sqrt D}
                  +\frac{S_*\sin^2x}{D^{3/2}}>0.
\end{equation}
Its branch-endpoint value is $1/\sin b$.
Compactness supplies a positive lower bound, retained on a thin
fixed convex rectangle and for $t$ in a small complex neighborhood.

On the remaining compact left arc, $\varphi=-\ii\gamma$ with
$\gamma>0$, so $v$ has real part $\theta$ and is injective.
There $1+g'(w)$ cannot vanish because $g'$ is purely imaginary.
Both endpoint slivers remain regular after a small extension.
The left image has real part less than $-b$, and the right image
has real part greater than $-b$; they meet in the regular branch
chart. An analytic inverse therefore exists along the embedded
compact arc, with bounded fixed jets after shrinking its neighborhood.
Choose the deformation and parameter margins inside this neighborhood.
\end{proof}

Let $\phi(t,v)=v-U(t,v)$ and
$D(t,v)=4\cos^2(v/2)-S(t)^2$.
The square root below is fixed by
$\sqrt D=\sin\phi-\sin U$; it is positive on the limiting right
arc and continued along the one-sided inverse. Implicit differentiation
of $\cos U+\cos\phi=S$ gives the exact formulas
\begin{equation}\label{eq:regular-explicit}
 U_v=\frac{\sin\phi}{\sqrt D},\qquad
 U_t=\frac{S'}{\sqrt D},\qquad
 a(t,v):=\frac{U_v}{\sinh(\ii\phi)}=-\frac{\ii}{\sqrt D}.
\end{equation}
The denominator never vanishes on the chosen neighborhoods.
These formulas directly prove that the regular one-body coefficient
and its joint jets are bounded, including at the inverse branch.

The ordinary bordered kernel is analytic in these regular coordinates.
Near the double branch, the even inverse in $\varphi$ cancels
$\varphi_i+\varphi_l$ from its denominator. Away from there the
denominator is separated: the double lower branch is the only limiting
zero of a pair denominator $1-z_i z_l$ on the one-sided support.
In particular,
\begin{equation}\label{eq:v-3}
 |L(v_i,v_l)|\le C|v_i-v_l|\quad\text{on }\mathcal D_R^2,
\end{equation}
and its fixed joint jets are bounded. The diagonal zero remains
in this single inverse coordinate even at the branch.

\subsection{Square-root geometry on the physical contour}
\begin{lemma}[Branch normal form]\label{lem:constrained-branch-geometry}
Put $u=\theta+b$ and $a_B=\sin b>0$.  On the original unit
contour, and on a true lower branch displaced outward by a common
$\mu_0\ge0$, set $E_0=\varepsilon+\mu_0$.  Uniformly for the fixed
small branch arc and $0\le \mu_0\le\tau$, one has
\[
 \begin{aligned}
 D&=1-W=a_Bu-i\beta+R,\qquad cE_0\le\beta\le CE_0,\\
 R&\in\mathbb R,\quad |R|\le C(u^2+E_0^2),\qquad
 D_u=a_B+O(|u|+E_0).
 \end{aligned}
\]
Here $\beta$ may depend smoothly on $u$, with
$|\beta_u|\le CE_0$.  On the physical fourth-quadrant branch,
\begin{equation}\label{eq:constrained-branch-bounds}
 \varphi=\sqrt{2D}\,h(D),\qquad h(0)=1,
 \qquad |\varphi'|\asymp(|u|+E_0)^{-1/2},
\end{equation}
where $h$ is analytic.  Moreover
\[
 -\Im\varphi\ge c\sqrt{|u|+E_0}\quad(u\le0),\qquad
 \Re\varphi'\ge\tfrac23|\varphi'|\quad(u\ge0),\qquad
 \int_{|u|\le\delta_B}|\varphi'(u)|\,du\le C.
\]
For coordinates with the same outward shift, a minimum nonnegative
coordinate $u_j$ satisfies $|\varphi'_i|\le C|\varphi'_j|$
for every other nonnegative coordinate.

\end{lemma}
\begin{proof}
Use $y=e^{\mu_0+i(-b+u)}$ and the exact real and imaginary parts of
$W=S-\cosh(\mu_0+i(-b+u))$.  The imaginary part is
\[
 \beta=\Im S+\sinh(\mu_0)\sin(b-u),
\]
and $\Im S=2\sin(\arg t_*)\varepsilon+O(\varepsilon^2)$.
The real part has the asserted expansion because
$1-S_*+\cos b=0$, $\Re S-S_*=O(\varepsilon^2)$, and
$\cosh \mu_0-1=O(\mu_0^2)$.  The analytic factorization of
$\arccos(1-D)$ gives \eqref{eq:constrained-branch-bounds}.
The relative error with respect to $a_Bu-i\beta$ is uniformly
small after the branch width and $\tau$ are fixed small.  Its
square root lies in a small fixed enlargement of the fourth quadrant.
For $u\le0$
its imaginary part has size at least a fixed multiple of its modulus;
for $u\ge0$ the derivative lies in a cone whose angle is at most
$\pi/4+o(1)$ from the positive real axis.  This proves the stated
attenuation and $2/3$ cone bounds.  The integrable majorant and the
minimum-coordinate comparison follow immediately from the two-sided
derivative estimate.

\end{proof}

\subsection{Coordinate size and angular diameter}

\begin{lemma}\label{lem:coordinate-scale}
On the actual right contour put
$u_i=\theta_i+b$, $\mu=N^{-1}\sum_i u_i$ and
$d=\max_i\theta_i-\min_i\theta_i>0$.
For any $r$ with $r\le d\le2r$, set $T=|\mu|+r+\varepsilon$.
Then
\begin{equation}\label{eq:v-21}
 |1+\varphi_w(w_i)|\le C(|u_i|+\varepsilon)^{-1/2},
\end{equation}
\begin{equation}\label{eq:v-22}
 \int_{\mu-2r}^{\mu+2r}(|u|+\varepsilon)^{-1/2}\dd u
       \le\frac{Cr}{\sqrt T},
\end{equation}
and
\begin{equation}\label{eq:v-23}
 |U_v(t,v_i)|\le C\sqrt T,\qquad
 |v_i-v_l|+|L_{il}|\le Cr/\sqrt T.
\end{equation}
These are pointwise estimates. No extrema, diameter, or mean scale
is differentiated, and no partition of the integral is made here.
\end{lemma}
\begin{proof}
Near the lower branch, $p=0$, $m=1$, so the outward displacement
is the same $\tau P/N$ at every such coordinate. The branch normal
form gives \eqref{eq:v-21}; away from that branch compactness gives
the same bound, with a larger constant. It also holds along a
one-variable angular curve with the actual occupancies held fixed.
The curve derivative has the additional bounded factor
$1-\ii\tau[(P/N)m'-(Q/N)p']=1+O(\tau)$.

If $|\mu|\ge4r$, the denominator in \eqref{eq:v-22} is comparable
to $\sqrt{|\mu|+r+\varepsilon}$. Otherwise enlarge the interval to
$[-6r,6r]$ and use
\[
 \sqrt{6r+\varepsilon}-\sqrt\varepsilon
 =\frac{6r}{\sqrt{6r+\varepsilon}+\sqrt\varepsilon}.
\]
This proves the interval estimate even when the interval crosses
the branch. Joining two actual angles along their fixed-occupancy
curve gives exactly the actual $v$ endpoints. Integrating its
derivative, and then using \eqref{eq:v-3}, proves the second
inequality in \eqref{eq:v-23}.

For the remaining estimate, first suppose $P=0$.
Nonnegativity of $p$ gives $h_i=0$ for every angle.
The branch normal form and compactness imply
$|\varphi_i|\le C\sqrt{|u_i|+\varepsilon}\le C\sqrt T$.
Equation~\eqref{eq:regular-explicit} then gives
$|U_v(t,v_i)|\le C\sqrt T$.
If $P>0$, some actual right angle lies in the positive support of
$p$, hence $\max u_i\ge b$ and $T\ge b/2$.
Boundedness of $U_v$ gives the same conclusion.
All estimates permit arbitrary actual occupancy fractions in $[0,1]$.

\end{proof}

\section{Tangential differentiation of the complete one-sided integral}\label{sec:constrained-transport}\label{sec:global-transport}\subsection{The scalar density and its full Jacobian}

Let $M=D_\theta w=I-\ii Dh$. Its inverse is uniformly bounded for
small $\tau$, $\det M$ has bound $C^N$, and fixed jets add polynomial
factors in $N$. Also $\boldsymbol1^TM=\boldsymbol1^T$, hence
$\boldsymbol1^TM^{-1}=\boldsymbol1^T$. For
$\chi_t(\theta)=(v_i)$,
\begin{equation}\label{eq:v-4}
D_\theta\chi_t=\operatorname{diag}(1/U_{v,i})M,
 \quad J=\det D_\theta\chi_t
       =\det M\prod_i(1/U_{v,i}).
\end{equation}
This identity is used at $\varepsilon>0$, where all factors are
nonzero. The apparent branch poles will be retained as integrable
Jacobian factors, never differentiated as unregularized amplitudes.

The complete deformed scalar density is $J f(t,v)$, where
\begin{equation}\label{eq:v-5}
f(t,v)=\frac1{N!(2\pi)^{N-1}}
 \frac{1+\sigma Z}{1-\sigma Z}\det\mathcal A\prod_i a(t,v_i).
\end{equation}
Formula \eqref{eq:v-5} includes the unchanged factor $2\pi$ converting the
ambient normalized density to constrained Haar normalization. The
entire occupancy Jacobian has cancelled from the regular amplitude.

For a residual field $B$ in $v$ coordinates define
\begin{equation}\label{eq:v-6}
V=-M^{-1}(\dot U+U_v B),
\end{equation}
where a dot denotes the fixed-$v$ derivative $\partial_t$ and products
inside the vector are componentwise. Differentiating
$U(t,\chi_t(\theta))=w(\theta)$ shows that \eqref{eq:v-6} is exactly
$(D_\theta\chi_t)^{-1}(\partial_t\chi_t-B)$.
If
\begin{equation}\label{eq:v-7}
\sum_i U_{v,i}B_i=-\sum_i\dot U_i,
\end{equation}
then $\sum_iV_i=0$ on an ambient neighborhood. If in addition
\begin{equation}\label{eq:v-8}
\sum_iB_i=0,
\end{equation}
then $(\partial_t-V\cdot\nabla_\theta)\sum_i\varphi_i=0$.
Indeed $\sum\varphi_i=\sum v_i-\sum w_i$, and
$\sum w_i=\sum\theta_i$ is frozen by tangency.

For every fixed smooth real angular weight $\omega$ the exact identity
is
\begin{equation}\label{eq:v-9}
\partial_t(\omega Jf)-\operatorname{div}_\theta(V\omega Jf)
 =\omega J[\partial_tf+\operatorname{div}_v(Bf)]
    -(V\cdot\nabla_\theta\omega)Jf.
\end{equation}
To verify the identity, put $A_\chi=D_\theta\chi_t$.
Commuting mixed derivatives of $\chi_t$ gives
\[
 (\operatorname{div}_vB)\circ\chi_t
 =\operatorname{tr}\!\left[A_\chi^{-1}
   \left(\partial_tA_\chi-\sum_iV_i\partial_iA_\chi\right)\right]
     -\operatorname{div}_\theta V.
\]
Jacobi's formula therefore gives
$J(\operatorname{div}_vB)\circ\chi_t
=\partial_tJ-\operatorname{div}_\theta(JV)$.
The product and chain rules prove \eqref{eq:v-9}. Because $\sum V_i=0$ identically, the intrinsic divergence on
$\sum\theta_i=c$, parametrized by the first $N-1$ angles, is
\[
 \sum_{i<N}(\partial_i-\partial_N)V_i=\sum_i\partial_iV_i.
\]
The same formula holds for $V$ multiplied by any scalar density.
Thus ambient divergence restricts to constrained divergence without
an additional measure factor. Integration by parts applies to the
real and imaginary parts of these complex fields.
The identity concerns scalar coefficients with respect to the original
real constrained measure on $\mathbb H_N^\sigma$. The bumps remain
smooth functions of the real angles. Formal $v$ differentiation is the
commuting complex frame $(D_\theta\chi_t)^{-1}\nabla_\theta$; its
coefficients and their derivatives are controlled below.

More explicitly, if $A_\chi=D_\theta\chi_t$, set
\[
 E_a=\sum_i(A_\chi^{-1})_{ia}\partial_{\theta_i},\qquad
 E_t=\partial_t-
       \sum_i(A_\chi^{-1}\partial_t\chi_t)_i\partial_{\theta_i}.
\]
Then $E_a\chi_b=\delta_{ab}$ and $E_t\chi_b=0$. Every commutator
$[E_a,E_b]$ or $[E_t,E_a]$ is a linear combination of the real
angular derivatives and annihilates all $\chi_b$. Invertibility of
$A_\chi$ therefore makes every commutator zero. This is a pointwise differential-operator calculation on the original
real domain, with the occupancies treated as smooth real-angle
functions.

\subsection{A bounded transport when a left root is present}\label{sec:left-gap}

Choose a single smooth right cutoff $r(\theta)$, equal to one for
$\theta+b\ge-\delta_B/2$ and zero for
$\theta+b\le-\delta_B$, with its harmless periodic return inside the
excluded upper plateau. Decompose the complete one-sided integral by
\[
 1=\prod_i r(\theta_i)+\left(1-\prod_i r(\theta_i)\right).
\]
Telescope the second term. Each of its at most $N$ terms has one fixed
named anchor $a$ in a compact left region, including on every cutoff
jet. That anchor has $|z_a|\le v_0<1$, and $|U_{v,a}|\ge c>0$.

On such a term take
\begin{equation}\label{eq:v-10}
B_a=-\frac{\sum_i\dot U_i}{U_{v,a}},\qquad B_i=0\ (i\ne a).
\end{equation}
Equation \eqref{eq:v-7} holds; equation \eqref{eq:v-8} is unnecessary because
$|1-\sigma Z|\ge1-v_0$. All fixed regular-coordinate and parameter jets of
$B$, $V$ and the regular amplitude have polynomial $N$ bounds, apart
from fixed one-body products costing $C_j^N$. In particular the
derivatives of $Z=\exp(-\ii\sum\phi_i)$ are bounded by a polynomial
in $N$, since $|Z|\le1$ at the actual contour and fixed jets of each
$\phi_i$ are bounded.

Hadamard applied to the complementary minors after fixed derivatives
therefore bounds the transported density by
\begin{equation}\label{eq:v-11}
C_j^N N^{-N/2+C_j}\prod_i |1+\varphi_w(w_i)|.
\end{equation}
The named compact factor is bounded. Eliminate its angle using the
constraint; every remaining factor has the independent majorant
$C(1+\operatorname{dist}(\theta,-b)^{-1/2})$. Integrating the product
costs $C^N$. This proves a uniform summable bound for every fixed
derivative of the whole left-root part, including all its mixed
branch configurations.

\subsection{One collective field on the entire right part}\label{sec:right-field}

Define the mixed covariance and angular sums
\[
 R=\sum_{i<l}(U_i-U_l)(U_{v,i}-U_{v,l}),\qquad
 G_i=\sum_l(U_i-U_l),\qquad A=-\sum_iU_{t,i}.
\]
Off full equality define the field by
\begin{equation}\label{eq:v-12}
 B_i=\frac{A G_i}{R}.
\end{equation}
The denominator is nonzero on the actual right contour, as proved
below. Pairing indices gives the exact identities
\begin{equation}\label{eq:v-13}
 \sum_iG_i=0,\quad \sum_i U_{v,i}G_i=R,\qquad
 \sum_iB_i=0,\quad \sum_iU_{v,i}B_i=-\sum_iU_{t,i}.
\end{equation}
Thus the same field freezes both required sums on the entire right
part. There are no divided pair differences in its formula.

\begin{lemma}[Protected mixed covariance]\label{lem:covariance}
On an actual right-contour diameter shell $r\le d\le2r$, put
$T=|\mu|+r+\varepsilon$ and $\Delta=r/\sqrt T$.
For fixed sufficiently small geometric choices,
\begin{equation}\label{eq:v-15}
 \Rea(e^{-\ii\pi/4}R)
 \ge\frac c{\sqrt T}\sum_{i<l}(\theta_i-\theta_l)^2
 \ge c\sqrt T\Delta^2.
\end{equation}
In particular the field is regular at every partial diagonal.
\end{lemma}
\begin{proof}
Hold the actual occupancies $\rho=P/N$, $\upsilon=Q/N$ fixed and
join any two angles along
\[
 w(x)=x-\ii\tau\{\rho m(x)-\upsilon p(x)\},\qquad
 v(x)=\varphi_t(w(x))+w(x).
\]
This auxiliary curve has exactly the actual endpoints and
$w'(x)=1+O(\tau)$. Its images lie in the fixed inverse domain
for every $0\le\rho,\upsilon\le1$ by the same small-displacement
bounds used in Lemma~\ref{lem:inverse}.

The positive real $U_{vv}$ on the limiting right interval has
arbitrarily small argument in its sufficiently thin neighborhood.
Near the lower branch, \eqref{eq:regular-explicit} and the
fourth-quadrant normal form put $U_v$ in an arbitrarily small
enlargement of the argument interval $[-\pi/2,0]$.
On the remaining compact right arc $U_v$ is close to a positive
real number. Consequently
\[
 \frac{\dd}{\dd x}U_v(t,v(x))
       =\frac{U_{vv}}{U_v}w'(x)
\]
lies in a fixed small enlargement of the cone $[0,\pi/2]$.
First fix the thin inverse-coordinate neighborhood and a
strict angular margin below $\pi$ for the limiting cone.
Shrink the branch width $\delta_B$ until its square-root
normal-form enlargement is smaller than that margin, and then
reduce $\tau$ and the parameter neighborhood. In particular,
the allowed bound on $\delta_B$ in (6.1) is only an upper bound:
$\delta_B$ may be chosen smaller independently of $N$.
The enlarged cone then has width less than $\pi$ with a fixed margin.
For $\theta_i>\theta_l$, its integral has the same cone, while
$U_i-U_l=w(\theta_i)-w(\theta_l)$ has argument $O(\tau)$
and modulus comparable to $\theta_i-\theta_l$.
Every product in $R$ therefore has a uniformly positive
$\pi/4$ projection.

Quantitatively, $|U_v|\le C\sqrt T$ along each joining curve.
To see this, its $v$-oscillation is at most $C r/\sqrt T=C\Delta$
by the fixed-occupancy interval estimate in
Lemma~\ref{lem:coordinate-scale}. The endpoint bound
$|U_v(t,v_i)|\le C\sqrt T$, bounded $U_{vv}$ on the convex
domain, and $\Delta\le\sqrt T$ give the assertion throughout
the curve. Since $|U_{vv}|\ge c$ and $|w'|\ge c$, integration
in the cone gives
\[
 |U_{v,i}-U_{v,l}|\ge c|\theta_i-\theta_l|/\sqrt T.
\]
The positive projection of each covariance product is thus at least
$c(\theta_i-\theta_l)^2/\sqrt T$.
Summing proves \eqref{eq:v-15}, including at all partial collisions.
\end{proof}

\begin{lemma}[Joint jets at the regular-coordinate scale]\label{lem:jets}
For every fixed total order $m=a+|\beta|$, on the same actual contour,
\begin{equation}\label{eq:v-14}
 |\partial_t^a\partial_v^\beta B_i|
 +|\partial_t^a\partial_v^\beta V_i|
 \le C_mN^{C_m}\Delta^{-1-m}.
\end{equation}
Only pointwise scale envelopes are used; no mean, diameter or scale is
differentiated. There is no separate inverse-$\varepsilon$ loss.
\end{lemma}
\begin{proof}
Write $D_{il}=U_i-U_l$ and $E_{il}=U_{v,i}-U_{v,l}$.
The actual contour and Lemma~\ref{lem:coordinate-scale} give
\[
 D_{il}=O(\sqrt T\Delta),\quad U_{v,i}=O(\sqrt T),\quad
 E_{il}=O(\Delta).
\]
Every positive fixed-$v$ parameter derivative of either difference
$D_{il}$ or $E_{il}$ is $O(\Delta)$, by bounded analytic divided
differences on the convex inverse domain. Their other fixed joint
jets are bounded; a first $v$ derivative of $D_{il}$ has the
better size $O(\sqrt T)$.

Leibniz' rule yields, for any joint derivative $\mathcal D^m$,
\begin{equation}\label{eq:sharp-R}
 \begin{aligned}
 |\mathcal D^mR|&\le C_mN^{C_m}\sqrt T\Delta^{2-m},&
 |\mathcal D^mR^{-1}|&\le C_mN^{C_m}T^{-1/2}\Delta^{-2-m},\\
 |\mathcal D^mG_i|&\le C_mN^{C_m}\sqrt T\Delta^{1-m}.&&
 \end{aligned}
\end{equation}
Here are the complete size checks. The $R$ jets at orders
$0,1,2$ are bounded by
$\sqrt T\Delta^2$, $\sqrt T\Delta$, and $\sqrt T$;
higher orders are bounded. At order one the possibly different
parameter term has size $O(\Delta^2)$, which is at most
$O(\sqrt T\Delta)$. At order two the possible sizes are
$\sqrt T$, $\Delta$, and $\Delta^2$, all bounded by
$C\sqrt T$. For $G_i$ the corresponding sizes at orders
$0,1$ are $\sqrt T\Delta$, $\sqrt T$, and higher jets are
bounded. Since $\Delta\le\sqrt T$ and maximum scales are
fixed, these tables imply the displayed envelopes at every fixed
order. The quotient rule and \eqref{eq:v-15} give the inverse
bound. All jets of $A$ are polynomially bounded in $N$.
Formula~\eqref{eq:v-12} now proves the estimate for $B$.

The smooth-frame coefficients are exactly
\[
 E_a\theta_i=(M^{-1})_{ia}U_{v,a},\qquad
 E_t\theta_i=(M^{-1}U_t)_i.
\]
Their fixed joint jets have polynomial $N$ bounds: differentiate
$M^{-1}$ using $D(M^{-1})=-M^{-1}(DM)M^{-1}$ and induct on
the bounded real angular derivatives of $M$ and analytic jets of $U$.
No $U_v^{-1}$ enters. This also bounds every smooth cutoff jet.
Formula~\eqref{eq:v-6} and the product rule give the same
estimate for $V$.
\end{proof}

At fixed $\varepsilon,N$, the complete product identity gives
collision degree $N(N-1)$ in the angular differences.
The field has one inverse diameter and each transport loses at most
two degrees. The full-equality puncture flux therefore vanishes
through order $j\le k$ whenever
\begin{equation}\label{eq:v-18}
 N^2-2j-3\ge4N_*+2>0,\qquad N\ge N_*+2.
\end{equation}
No partial-diagonal flux occurs, since Lemma~\ref{lem:covariance}
separates the sole field denominator at every nonconstant configuration. The transports are performed on the
entire smooth right term before any diameter division.

There are two useful bounds for the resulting density. First,
the complete pair product has the exact regular factorization
\[
 P_{il}=(v_i-v_l)^2A_{il}(t,v_i,v_l),
\]
where $A_{il}$ and its fixed jets are analytic and bounded on
$\mathcal D_R^2$. This follows from \eqref{eq:kernel-square}
and the analytic diagonal zero in \eqref{eq:v-3}, including at
the double branch. Thus its $N(N-1)$ collision degrees remain
available.

For clarity, fix the derivative budget $j\le k$ and induct on the
number $r\le j$ of transports
$\partial_t+\operatorname{div}_v(B\,\cdot)$, including the
cutoff terms in \eqref{eq:v-9}. At $r=0$ the complete pair product
has the displayed $N(N-1)$ zeros. Under one further transport,
a spatial derivative can lose one collision degree and the
corresponding field coefficient at most one more;
$\operatorname{div}_v B$ has the same two-degree budget.
Parameter derivatives of the analytic pair factors preserve their
diagonal zeros, while derivatives of already formed field
coefficients satisfy \eqref{eq:v-14}. The product rule therefore
reduces the remaining collision degree by at most two per
transport. There are $O(N^2)$ choices of a differentiated pair
at each fixed step; the product-rule multiplicity over $r\le j$
is at most $N^{O(j)}$. The fixed-order jets of the smooth weights
and of $\det M$ obey the same polynomial budget by
Lemma~\ref{lem:jets}.

The transport freezes $\sum_i\varphi_i$ exactly by
\eqref{eq:v-7}--\eqref{eq:v-8}, so no derivative of
$(1-\sigma Z)^{-1}$ occurs. The remaining one-body factors
are retained in the right-hand side, and all analytic pair
quotients on the fixed inverse domain cost at most $C_j^{N^2}$.
Induction gives
\begin{equation}\label{eq:full-product-density}
 C_j^{N^2}\Delta^{N(N-1)-2j}
       \frac{\prod_i|1+\varphi_w(w_i)|}{|1-\sigma Z|}.
\end{equation}
Since $j$ is fixed, the polynomial factors
$N^{O(j)}$ and the one-body, cutoff and Jacobian costs
$C_j^N N^{C_j}$ can be absorbed into $C_j^{N^2}$,
uniformly for all odd $N\ge N_*+2$.
The factorial $1/N!$ can only improve the bound.

Second, on any region $d\ge d_0>0$, all field jets are bounded
with polynomial $N$ costs. Differentiate the bordered determinant
and apply Lemma~\ref{lem:border-minors} to its complementary
ordinary minors. Division by $N!\ge(N/e)^N$ gives
\begin{equation}\label{eq:far-density}
 C_{j,d_0}^N N^{-N/2+C_j}
       \frac{\prod_i|1+\varphi_w(w_i)|}{|1-\sigma Z|}.
\end{equation}
Both estimates apply to the same transported density; no cutoff
between them has been differentiated.

\section{One coarea estimate for the complete right integral}\label{sec:global-integration}
\subsection{A projected periodic kernel}
\begin{lemma}[Oblique periodic kernel]\label{lem:oblique-kernel}
Fix $|\vartheta|<\pi/2$ and put $P_\vartheta=2\pi\cos\vartheta$.
If $G=x-ia$, $a\ge a_0$, and $0<a_0\le1$, then, with
$\nu=\Re(e^{-i\vartheta}G)$,
\[
 \frac1{|1-e^{-iG}|}
 \le \frac{C_\vartheta}
 {a_0+\operatorname{dist}(\nu,P_\vartheta\mathbb Z)}.
\]
Consequently, on any injective $\nu$-parametrized arc whose
$\nu$-image has length $L$, the integral of this majorant in $\nu$
is at most $C_\vartheta(1+L)(1+|\log a_0|)$.
\end{lemma}
\begin{proof}
For $0<a\le1$, use
\[
 |1-e^{-a-ix}|^2=(1-e^{-a})^2+2e^{-a}(1-\cos x)
 \ge c\{a^2+\operatorname{dist}(x,2\pi\mathbb Z)^2\}.
\]
It follows that
$|1-e^{-iG}|\ge c\{a+\operatorname{dist}(x,2\pi\mathbb Z)\}$.
For an integer $n$ minimizing the latter distance,
\[
 |\nu-2\pi n\cos\vartheta|
 \le\cos\vartheta\,|x-2\pi n|+|\sin\vartheta|a.
\]
Thus $a_0+\operatorname{dist}(\nu,P_\vartheta\mathbb Z)$ is
bounded by a constant times the preceding lower bound.
For $a\ge1$, $|1-e^{-iG}|\ge1-e^{-1}$ and the proposed
denominator is bounded above by $1+P_\vartheta/2$.
Integration over translates of one period proves the final claim.
\end{proof}


For $\sigma=-1$, replace $G$ by $G-\pi$ in this lemma;
$1+e^{-\ii G}=1-e^{-\ii(G-\pi)}$. Only a fixed translate of
the projected period lattice changes. The derivative, interval length,
and all bounds remain the same.

For the actual contours the attenuation floor is
$-\log|Z|\ge cN\varepsilon$. Indeed $\Im W_i\ge c\varepsilon$,
and $W_i=\cos\varphi_i$ remains in a bounded set. Writing
$\varphi_i=x_i-ia_i$ gives
$\Im W_i=\sin x_i\sinh a_i\le C a_i$, hence
$a_i\ge c\varepsilon$. Thus Lemma~\ref{lem:oblique-kernel}
costs at most $C(1+L)(1+|\log\varepsilon|)$ on an injective
projected phase interval of length $L$.

\subsection{Oblique extreme coarea with the full occupancy dependence}

Let $I=[-b-\delta_B,\delta_R]$, where $\delta_R=\delta<b/2$,
and choose the negative segment inside the lower branch plateau.
The upper inverse branch is outside $I$.  Use the original balanced
contour with reflected disjoint bumps and put
\[
 \rho=P/N,\quad\upsilon=Q/N,\quad
 f(u;\rho,\upsilon)
 =-b+u-i\tau\{\rho m(-b+u)-\upsilon p(-b+u)\}.
\]
Here $u\in[-\delta_B,b+\delta_R]$ and the original one-sided
occupancies satisfy $P\le1/2$, $0\le Q\le N$.
The auxiliary analytic estimates in terms of $E=\varepsilon+\tau\rho$
hold on the fixed box $0\le\rho,\upsilon\le1$.
The final bound in Lemma~\ref{lem:global-coarea} concerns actual
occupancies computed from the angles, including on the enlarged
rectangle used below.
Let $\varphi_t(w)=\arccos(S(t)-\cos w)$ on the physical sheet,
$b_t(w)=\partial_w\varphi_t(w)$. Use the regular coordinate
defined above.
The regular phase coordinate is
$v_t(w)=\varphi_t(w)+w$.
Its analytic inverse and the scalar transport were proved in
Sections~\ref{sec:global-coordinate}--\ref{sec:global-transport}.

\begin{lemma}[Occupancy-dependent extreme derivative]\label{lem:global-coarea}
Write
$G(\theta)=\sum_i\varphi_t(w_i(\theta))$ and
$\vartheta=3\pi/8$.  On an actual extreme-index cell with
$0\le m=\min_i u_i<M=\max_i u_i$, all other $u_i$ fixed,
and $m+M$ fixed by the constraint, one has
\begin{equation}\label{eq:g-G1}
\Re\left(e^{-i\vartheta}\frac{dG}{dm}\right)>0,
 \qquad
 \frac{|(b_t(w_m)+1)(b_t(w_M)+1)|}
 {\Re(e^{-i\vartheta}dG/dm)}
 \le C\frac{\sqrt{M+\varepsilon}}{M-m}.
\end{equation}
The constants are independent of $N$, $\varepsilon$, and the
actual occupancies.  A sufficiently small fixed $\tau$ is chosen
after the bump derivatives and fixed angular margins.
\end{lemma}
\begin{proof}
Temporarily hold the actual values of $\rho,\upsilon$ fixed
and set $E=\varepsilon+\tau\rho$.
On the plateau $f=-b+u-i\tau\rho$; the square-root normal form gives
\begin{equation}\label{eq:g-G2}
|b_t(f(u))|\le C(u+E)^{-1/2},\qquad
 |\varphi_t''(f(u))|\asymp(u+E)^{-3/2},\qquad u\ge0.
\end{equation}
These estimates extend to the remainder of the fixed right
interval by compactness: it stays a fixed distance from both
dispersion branches, while $u$ is bounded above and below there.
The existing inverse supplies the curvature identity
\[
 -\varphi_t''(w)=\frac{U_{vv}(t,v)}{U_v(t,v)^3},
 \qquad v=\varphi_t(w)+w,
\]
obtained by differentiating $\varphi_w=U_v^{-1}-1$.
On the true branch with $u\ge0$, the normal form puts $U_v$
in an arbitrarily small enlargement of $[-\pi/4,0]$ in argument.
On the compact remainder $U_v$ is close to a positive real number.
The argument of $U_{vv}$ is uniformly small, as in
Lemma~\ref{lem:covariance}. Hence, for $\vartheta=3\pi/8$,
\begin{equation}\label{eq:g-G3}
 \Re\{e^{-\ii\vartheta}[-\varphi_t''(f(u))]\}
       \ge c|\varphi_t''(f(u))|.
\end{equation}
Indeed the curvature arguments lie in a small enlargement of
$[0,3\pi/4]$, whose midpoint is $3\pi/8$.
The same identity and the branch size of $U_v$ give the
second bound in \eqref{eq:g-G2} directly.

Put $a(u)=\partial_u f(u)$ and $B(u)=b_t(f(u))a(u)$ at these
fixed occupancies.  Since $a=1+O(\tau)$,
$a'=O(\tau)$, and $|b_t(f)|\le C|\varphi_t''(f)|$ by \eqref{eq:g-G2},
\[
 -B'=-\varphi_t''(f)a^2-b_t(f)a'
\]
retains \eqref{eq:g-G3}.  Therefore, with $d=M-m$,
\begin{equation}\label{eq:g-G4}
\begin{aligned}
 &\Re\{e^{-i\vartheta}[B(m)-B(M)]\}\ge cL,\\
 &L=\int_m^M(u+E)^{-3/2}du
 =\frac{2d}{\sqrt{m+E}\sqrt{M+E}
 (\sqrt{m+E}+\sqrt{M+E})}.
\end{aligned}
\end{equation}

We now restore the entire occupancy dependence rather than
moving the occupancies along that auxiliary comparison curve.
Set
\[
 B_m^{\rm occ}=\sum_v m(\theta_v)b_t(w_v),\qquad
 B_p^{\rm occ}=\sum_v p(\theta_v)b_t(w_v).
\]
Direct differentiation of the original balanced contour gives
the exact formula
\begin{equation}\label{eq:g-G5}
\partial_{\theta_i}G
 =b_t(w_i)a(u_i)
 -\frac{i\tau p'(\theta_i)}N B_m^{\rm occ}
 +\frac{i\tau m'(\theta_i)}N B_p^{\rm occ}.
\end{equation}
The constraint slice differentiates in the direction
$\partial_{\theta_{\min}}-\partial_{\theta_{\max}}$.
Every actual $u_v\ge m$, and the true branches have the same
outward displacement $\tau\rho$.  Hence \eqref{eq:g-G2} implies
\[
 \frac{|B_m^{\rm occ}|+|B_p^{\rm occ}|}{N}
 \le C(m+E)^{-1/2}.
\]
Both bump-derivative differences in \eqref{eq:g-G5} are bounded by $Cd$.
The total occupancy correction to $B(m)-B(M)$ is thus at most
$C\tau d(m+E)^{-1/2}$.
Since $M+E$ lies in a fixed bounded interval, the exact
expression in \eqref{eq:g-G4} gives
$d(m+E)^{-1/2}\le CL$.  Choosing $\tau$ small absorbs the
correction into half the positive projection in \eqref{eq:g-G4}.
This comparison uses the actual extreme configuration only;
it does not presume a Hessian bound along a transposition path
on which unchanged indices might become more extreme.

The numerator in \eqref{eq:g-G1} is at most
$C[(m+E)(M+E)]^{-1/2}$, with the constant enlarged on the bounded interval.
Combining with \eqref{eq:g-G4} bounds the ratio by $C\sqrt{M+E}/d$.
If $P=0$, then $E=\varepsilon$.  If $P>0$, an actual angle
belongs to the positive support of $p$, so $M\ge b$.
Therefore $E\le C(M+\varepsilon)$ in either event, proving
\eqref{eq:g-G1}.
\end{proof}



\paragraph{The geometric integration domain.}\label{par:right-rectangle}
After all differentiation, bound the smooth weight and its
transported coefficients by their proved envelopes and enlarge
the nonnegative integral to the full right angular rectangle
$I^N\cap\mathbb H_N^\sigma$. All actual occupancies there still satisfy
$0\le\rho,\upsilon\le1$, so the preceding lemma applies.
The balanced displacement has the same sign on each semicircle,
and hence $\Im(S-\cos w_i)\ge\Im S$: the physical sheet and
attenuation floor are preserved on this enlarged domain.
The regular inverse also remains valid there. On the compact
right arc its $v$-image changes by $O(\varepsilon+\tau)$;
near the branch, uniform square-root continuity bounds the
change by $O(\sqrt{\varepsilon+\tau})$. Choose $\delta_B$,
$\tau$ and $\varepsilon_0$ after the fixed margin of
$\mathcal D_R$ so that all these images lie inside it.

For fixed other coordinates and a fixed constrained sum,
the conditions specifying the extreme indices, a diameter
shell, and a nonnegative minimum are linear interval conditions
in the remaining variable $m$, with $M=A-m$.
Their intersection is therefore one interval. The negative
part is bounded directly below. No assertion about the number
of components of a smooth-cutoff fiber is required.

\subsection{A fixed mean scale for every diameter}


Write $u_i=\theta_i+b$, $\mu=N^{-1}\sum_i u_i$, and choose actual
minimum $m$ and maximum $M$ only in the final absolute integral.
The angular constraint fixes $\mu=b+(2\pi l+a_\sigma)/N$; at most $CN$ such means
meet the right interval. The independent coordinate measure after
eliminating any one angle is unchanged.

On a shell $r\le d\le2r$ use the scale
$T=|\mu|+r+\varepsilon$ from Lemma~\ref{lem:coordinate-scale}.
Every coordinate belongs to $[\mu-2r,\mu+2r]$.

For a nonnegative minimum, Lemma~\ref{lem:global-coarea}
gives the monotone projected phase and
\begin{equation}\label{eq:v-24}
 \frac{|1+\varphi_w(w_m)|\,|1+\varphi_w(w_M)|}
 {|d\nu/dm|}\le C\frac{\sqrt{M+\varepsilon}}{M-m}.
\end{equation}
Its phase range is at most $CN$, since each phase is bounded.
Lemma~\ref{lem:oblique-kernel} therefore costs $CN(H+1)$.
Here and below $H=|\log\varepsilon|$.

For a nonnegative minimum, \eqref{eq:v-24} therefore bounds the integral in
the remaining extreme coordinate by
\begin{equation}\label{eq:v-25}
\int \frac{|1+\varphi_w(w_m)|\,|1+\varphi_w(w_M)|}{|1-\sigma Z|}\dd m
       \le CN\frac{\sqrt T}{r}(H+1).
\end{equation}
For a negative minimum, $m$ lies in the fixed true-branch segment,
where its attenuation gives
\[
 \frac{|1+\varphi_w(w_m)|\,|1+\varphi_w(w_M)|}{|1-\sigma Z|}
 \le \frac{C}{(|m|+\varepsilon)\sqrt{|M|+\varepsilon}}.
\]
Set $a=|m|+\varepsilon$, $b_1=|M|+\varepsilon$. Since
$a+b_1\ge M-m\ge r$, the elementary inequality
\begin{equation}\label{eq:v-26}
\frac1{a\sqrt{b_1}}
 \le \frac1{\sqrt r}\left(\frac1a+\frac1{b_1}\right)
\end{equation}
follows by multiplying through by $a\sqrt{b_1}$ and using
$a+b_1\ge r$. On the constrained slice $M=A-m$, so each reciprocal
in \eqref{eq:v-26} integrates to $C(H+1)$ over a bounded interval. This gives
\begin{equation}\label{eq:v-27}
\int \frac{|1+\varphi_w(w_m)|\,|1+\varphi_w(w_M)|}{|1-\sigma Z|}\dd m
 \le Cr^{-1/2}(H+1)\le C\frac{\sqrt T}{r}(H+1),
\end{equation}
which is the same pair envelope as \eqref{eq:v-25}. 

The two actual extreme Jacobian factors have been retained through
\eqref{eq:v-25}; only the $N-2$ remaining factors are now replaced by
independent one-body majorants. Their interval integrations cost
$(Cr/\sqrt T)^{N-2}$. Thus the full coarea and interval cost is
\[
 C^NN^{C}\Delta^{N-3}(H+1),\qquad \Delta=r/\sqrt T.
\]
This includes the polynomial number of extreme-index cells and means.
Combined with \eqref{eq:full-product-density}, the complete shell is
bounded by
\begin{equation}\label{eq:v-28}
 C_j^{N^2}\Delta^{N^2-2j-3}(H+1).
\end{equation}
For every tail order $N\ge N_*+2$ and $j\le k$, the exponent
is at least $4N_*+2>0$.

Choose a fixed $\delta_0>0$ after all constants for $j\le k$,
small enough that $C_j\delta_0\le1/2$ for each of them.
On shells with $\Delta\le\delta_0$, the dyadic sum converges
uniformly: for fixed $\mu,\varepsilon$,
\[
 \frac{\Delta(r/2)}{\Delta(r)}
 =\frac12\sqrt{\frac{|\mu|+r+\varepsilon}
 {|\mu|+r/2+\varepsilon}}\in[1/2,1/\sqrt2].
\]
Its largest term is bounded by
\[
 C_j^{N^2}\delta_0^{N^2-2j-3}
 \le \delta_0^{-2k-3}2^{-N^2},
\]
which is summable over all odd tail orders.

On the complementary shells, $\Delta\ge\delta_0$ and
$T\ge r$ imply $r\ge\delta_0^2$. There are only finitely many
such dyadic angular shells, uniformly in $N,\varepsilon$.
Use \eqref{eq:far-density}; the same coarea and interval
integration is at most $C_j^NN^{C_j}(H+1)$, since
$\delta_0\le\Delta\le C$.
This gives the summable envelope
\[
 C_j^NN^{-N/2+C_j}(H+1).
\]
Both parts of this right-side sum include every constrained mean and
every odd tail order $N\ge N_*+2$.

\section{The complete logarithmic tail}
\label{sec:constrained-tail-sum}

\begin{theorem}\label{thm:constrained-tail}
At each positive-geometry signed point and for every $0\le j\le k$, uniformly
for $0<\varepsilon<\varepsilon_0$,
\begin{equation}\label{eq:logarithmic-absolute-tail}
 \sum_{\substack{N>N_*\\N\ \mathrm{odd}}}
    |\mathcal C_N^{(j)}((1+\varepsilon)t_*)|
       \le C_j(1+|\log\varepsilon|).
\end{equation}
In particular the entire higher-particle contribution is
$o(\varepsilon^{-1})$ at the derivative order of the first
singularity.
\end{theorem}
\begin{proof}
The selected part has the summable bound proved in
Lemma~\ref{lem:constrained-selected}, and the left-root part has
the bound \eqref{eq:v-11}. For the entire right part,
\eqref{eq:v-28} and the following small/large $\Delta$ division
give a summable majorant
\[
 \left(C_j2^{-N^2}+C_j^NN^{-N/2+C_j}\right)(H+1).
\]
These estimates cover the exact identity $\mathcal C_N=F_N+2L_N$.
Summing proves the theorem, and
$\varepsilon(1+|\log\varepsilon|)\to0$ gives the required
noncancellation estimate.

The choices have the following order: the selected point and $k$;
the endpoint slivers and fixed inverse-coordinate neighborhood;
the branch width, smooth bumps and right cutoff; a sufficiently
small fixed $\tau$; the parameter-disk constants; and then
$\varepsilon_0$ and the fixed estimate threshold $\delta_0$.
The branch width may also meet the finite gap required only for
Section~\ref{sec:first-single}.
All choices are independent of the tail order $N$.
The division by $\Delta$ is made only after the complete weighted
integrals have been transported, so creates no extra derivative
or boundary term.
\end{proof}


\section{The first singularity on the constrained angular torus}
\label{sec:first-single}
Fix a prime-family point $\zeta=e^{\ii\alpha}$,
$0<\alpha<\pi$, $N=p$ and $k=(p^2-1)/2$.
We first justify its finite-order complement using the same
signed positive-geometry construction as for the tail.

\begin{lemma}[Finite-order complement]\label{lem:finite-complement}
Every fixed derivative of every odd coefficient $n<p$ is
$O(1+|\log\varepsilon|)$ along $(1+\varepsilon)\zeta$.
At $n=p$, remove two small fixed reflected angular charts around
$\theta_i=-\alpha$ and $\theta_i=+\alpha$ for all $i$.
Every fixed derivative of the remaining complete coefficient is
$O(1+|\log\varepsilon|)$ on this radial approach.
\end{lemma}
\begin{proof}
The case $n=1$ follows from $C_1=1/(S-2)$.
For negative $S_*$ use \eqref{eq:negative-derivatives}; thus
it suffices to work at positive $S_*$ with the appropriate
constraint $\sigma$ and angle $\widehat\alpha\in(0,\pi/2)$.
The signed $n=1$ formula is $C_1^-=1/(S+2)$.
All remaining orders satisfy $n\ge3$.

The finite mean gap in \eqref{eq:finite-mean-exclusion}
places all constrained equality means away from the branch.
For each fixed $n$, the mixed-population part $F_n$ has all
true branches displaced by at least $\tau/(4n)$ and a fixed
product damping gap. For each fixed $n$, compactness gives a fixed
parameter disk on which its density is holomorphic and bounded.
This is a disk for the analytic continuation of the deformed
integrand from the physical parameter: the displacement keeps
the true branches separated from $W=\pm1$, the other roots
stay in compact regular neighborhoods, and the product damping
gap remains positive after shrinking the disk. This reasoning
also applies when the disk crosses $\Im S(t)=0$;
the analytic continuation, rather than a new pointwise
physical-sheet choice, is used in Cauchy's estimate.
Every left-root term has bounded fixed derivatives by the
anchor field of Section~\ref{sec:left-gap}, with its compact
angle eliminated before integrating the remaining one-body
majorants.

At full equality the constrained mean is
$\theta_0=(2\pi\ell+a_\sigma)/n$.
The balanced shift is exactly zero there, since
$P=np(\theta_0)$ and $Q=nm(\theta_0)$.
The product denominator vanishes only when
\[
 n\arccos(S_*-\cos\theta_0)\in2\pi\mathbb Z+a_\sigma.
\]
For $\sigma=1$ this is the ordinary Nickel equation; for
$\sigma=-1$ translate both angles by $\pi$ and use oddness of
$n$. The arithmetic of Section~\ref{sec:prime-roots} excludes
it at $n<p$, and at $p$ leaves only the lower center
$-\widehat\alpha$ on the one-sided support.
All other equality neighborhoods have regular roots and a
separated product denominator, so their localized integrals
extend holomorphically through the point.

Remove these neighborhoods and, at $p$, the critical chart,
chosen entirely inside $m=1,p=0$. There the balanced deformation
is the identity, $w_L=1$, and the right cutoff is one.
On the remaining right support, including all cutoff jets,
$d\ge d_0>0$. The collective field is therefore bounded to
every fixed jet order and has no puncture flux.
Transport gives a bound
\[
 C_{n,j}\frac{\prod_i|1+\varphi_w(w_i)|}{|1-\sigma Z|}.
\]
Retain the two extreme-coordinate factors for the coarea
estimate before bounding the others by independent majorants.
There are finitely many constrained means and shells with
$d\ge d_0$, so Section~\ref{sec:global-integration} gives
$O(1+|\log\varepsilon|)$ for every fixed derivative of this
remainder. Smooth-cutoff fibers are enlarged to the same
geometric intervals used there.

Reflection in the complete identity $\mathcal C_n=F_n+2L_n$
accounts for the two removed charts. For the negative half, negation and conjugation map
these charts back to $\pm\alpha$ and preserve the bound.
\end{proof}

Near the lower physical chart write
$\theta_i=-\alpha+u_i$, $u_N=-\sum_{i<N}u_i$.
Because $p\alpha\in2\pi\mathbb Z$, this is the literal
constraint with measure $\prod_{i<N}\dd u_i/(2\pi)$.
At $t_\varepsilon=(1+\varepsilon)\zeta$, the physical phase satisfies
\[
 \varphi_i=\alpha-2\ii\varepsilon+u_i-\cot\alpha\,u_i^2
      +O(\varepsilon^2+\varepsilon|u_i|+|u_i|^3).
\]
Indeed $\varphi_t=-2\ii\zeta^{-1}$,
$\varphi_u=1$ and $\varphi_{uu}=-2\cot\alpha$ at the center.
For $q(t,u)=1-Z(t,u)$ this gives
\begin{equation}\label{eq:single-q}
 q(t_\varepsilon,u)=2p\varepsilon-\ii d_\alpha\sum_i u_i^2
       +O(\varepsilon^2+\varepsilon|u|+|u|^3),\quad
 d_\alpha=\cot\alpha\ne0,\quad
 q_t(\zeta,0)=2p\zeta^{-1}=:D_t.
\end{equation}
The two leading terms have perpendicular phases. Consequently,
on a sufficiently small fixed shape ball,
\begin{equation}\label{eq:single-q-bound}
 |q(t_\varepsilon,u)|\ge c(\varepsilon+|u|^2).
\end{equation}
To check the remainder, bound it by
$C(\varepsilon_0+\delta)(\varepsilon+|u|^2)$ for
$|u|\le\delta$, $\varepsilon\le\varepsilon_0$, and choose
both radii so this is less than half the leading lower bound.
This argument works for either sign of $d_\alpha$.

With a fixed smooth shape cutoff $\chi$ equal to one near zero,
the lower chart is exactly
\begin{equation}\label{eq:angular-local-integral}
 I_p(t)=\int_{\mathbb R^{p-1}}
       \frac{\chi(u)\Delta(u)^2 A(t,u)}{q(t,u)}
                         \,\dd u_1\cdots\dd u_{p-1}.
\end{equation}
Here $\Delta(u)=\prod_{i<j}(u_i-u_j)$ and $A$ is analytic.
Its value, including all physical factors, is
\begin{equation}\label{eq:angular-local-prefactor}
 A(\zeta,0)=-\ii a_p,\qquad
 a_p=\frac{2^{1-p(p-1)}}{p!(2\pi)^{p-1}\sin^{p^2}\alpha}>0.
\end{equation}
Each sine numerator gives half the angle difference,
each hyperbolic denominator tends to $\ii\sin\alpha$,
each one-body factor tends to $(\ii\sin\alpha)^{-1}$,
and $1+Z\to2$. Their combined phase is
$(-\ii)^p(-1)^{p(p-1)/2}=\ii^{-p^2}=-\ii$.
Every parameter derivative retains the full Vandermonde
because the real shape coordinates are fixed.
The upper chart is exactly equal under reflection.
There are two centers; exchanging equal dispersion angles creates
no additional center.

\begin{lemma}[A nonzero radial period]\label{lem:single-period}
Let
\[
 \mathcal S_p=p^{-1/2}
       \int_{\substack{\sum_i\omega_i=0\\|\omega|=1}}
            \Delta(\omega)^2\,\dd\sigma(\omega)>0.
\]
Then, with $v_p=-\sum_{i<p}v_i$,
\begin{equation}\label{eq:single-period}
 J_\alpha=\int_{\mathbb R^{p-1}}
       \frac{\Delta(v)^2}{(2p-\ii d_\alpha\sum_i v_i^2)^{k+1}}
                \,\dd v_1\cdots\dd v_{p-1}
       =\frac{\mathcal S_p}{2k(2p)|d_\alpha|^k}>0.
\end{equation}
\end{lemma}
\begin{proof}
The coordinate measure on the zero-sum hyperplane is its
Euclidean measure divided by $\sqrt p$.
The Vandermonde is positive on an open spherical set after
squaring, so $\mathcal S_p>0$.
The radial exponent is $p^2-2=2k-1$.
With $Q=2p$ and $c=-\ii d_\alpha$, put $x=|v|^2$ and use
\[
 \frac{\dd}{\dd x}\left(\frac{cx}{Q+cx}\right)^k
       =kQc^k\frac{x^{k-1}}{(Q+cx)^{k+1}}.
\]
The bracket has endpoint values zero and one, and $Q+cx$ never
vanishes for $x\ge0$. The integral is absolutely convergent
with radial tail $R^{-3}$.
Since $k\equiv0\pmod4$, $c^k=|d_\alpha|^k$.
Polar coordinates now give \eqref{eq:single-period} directly.
\end{proof}

\begin{theorem}[The complete first odd singularity]\label{thm:single-first}
At the selected prime root,
\[
 C_p^{(k)}(t_\varepsilon)=\Lambda_\alpha\varepsilon^{-1}
           +o(\varepsilon^{-1}),\qquad
 \Lambda_\alpha=-2\ii a_p k!(2p)^k\zeta^{-k}J_\alpha\ne0.
\]
\end{theorem}
\begin{proof}
The complement is $O(1+|\log\varepsilon|)$ by
Lemma~\ref{lem:finite-complement}.
Differentiate \eqref{eq:angular-local-integral} $k$ times.
The unique highest pole is
\[
 \int \frac{\chi(u)\Delta(u)^2(-1)^k k!
       q_t(t_\varepsilon,u)^k A(t_\varepsilon,u)}
                {q(t_\varepsilon,u)^{k+1}}\,\dd u.
\]
Scale $u=\sqrt\varepsilon\,v$. The exponent is
$[p(p-1)+(p-1)]/2-(k+1)=-1$.
Equation~\eqref{eq:single-q-bound} gives the integrable majorant
$C\Delta(v)^2/(1+|v|^2)^{k+1}$.
Dominated convergence gives the single-chart coefficient
$-\ii a_p k!(2p)^k\zeta^{-k}J_\alpha$.
Every lower pole power $m+1$, $m<k$, contributes $o(\varepsilon^{-1})$,
since
\[
 \varepsilon\int_0^\delta
       \frac{r^{2k-1}}{(\varepsilon+r^2)^{m+1}}\,\dd r\longrightarrow0.
\]
At $m=k-1$ this is $O(\varepsilon|\log\varepsilon|)$;
at smaller $m$ it is $O(\varepsilon)$.
The reflected chart gives the identical coefficient, proving the
factor two. All phases and estimates hold for both signs of
$\cos\alpha$.
\end{proof}

\section{Noncancellation and the whole natural boundary}
Let
\[
 V(t)=\sum_{\substack{N\ge1\\N\ {\rm odd}}}C_N(t).
\]
Normal convergence permits termwise differentiation in $|t|>1$.
At a selected $\zeta=e^{\ii\alpha}$, the finitely many lower
orders are logarithmically bounded by Lemma~\ref{lem:finite-complement},
the first term is given by Theorem~\ref{thm:single-first},
and the entire higher odd tail is logarithmically bounded by
Theorem~\ref{thm:constrained-tail}, using
\eqref{eq:negative-derivatives} on the second quadrant.
Therefore
\[
 V^{(k)}((1+\varepsilon)\zeta)
       =\Lambda_\alpha\varepsilon^{-1}+o(\varepsilon^{-1}),
 \qquad \Lambda_\alpha\ne0.
\]
A holomorphic continuation through $\zeta$ would bound this
finite derivative on a short radial segment, which is impossible.
The factor $Q_H$ is locally analytic and nonzero at every selected
point. If $\X(1/t)=Q_H(t)V(t)$ continued there, division by
$Q_H$ would continue $V$. The same argument applies to
$\chi_H(1/t)=\beta_{\rm phys}(1/t)Q_H(t)V(t)$, because the
physical branch of $\beta_{\rm phys}(s)=\operatorname{arsinh}(s)/(2J)$
is locally analytic and nonzero at every selected nonaxis point.

The selected primes are dense in both upper quadrants.
Complex conjugation supplies a dense obstruction in the lower
quadrants. Every neighborhood of every unit-circle point,
including all four axis points, contains an obstructing point.
Thus no point of that circle admits local holomorphic continuation
of the exterior germ. The map $s=1/t$ is biholomorphic near every
unit-circle point and gives the assertion of
Theorem~\ref{thm:main} for the prescribed high-temperature germ.


\begin{thebibliography}{9}
\bibitem{Orrick}
W. P. Orrick, B. G. Nickel, A. J. Guttmann, and J. H. H. Perk,
The susceptibility of the square lattice Ising model: New developments,
\emph{Journal of Statistical Physics} \textbf{102} (2001), 795--841.
\href{https://arxiv.org/abs/cond-mat/0103074}{arXiv:cond-mat/0103074}.
\bibitem{Boukraa}
S. Boukraa, A. J. Guttmann, S. Hassani, I. Jensen, J.-M. Maillard,
B. Nickel, and N. Zenine, Experimental mathematics on the magnetic
susceptibility of the square lattice Ising model,
\emph{Journal of Physics A: Mathematical and Theoretical}
\textbf{41} (2008), 455202.
\href{https://arxiv.org/abs/0808.0763}{arXiv:0808.0763}.
\end{thebibliography}
\end{document}
